\documentclass{situationist-preprint}

\title{The Rate of Passage Objection}
\msid{TSM-02-b}
\msfamily{TSM-02}
\msdiscipline{Time, Space and Modality}
\mstype{Companion manuscript}
\mskeywords{passage of time; rate of passage; one second per second; hypertime; B-facts}

\begin{document}
\maketitle

\begin{abstract}
If time passes, it seems that it must pass at some rate, and the only candidate, one second per second, has been said to be no rate at all. Defenders of passage reply that a dimensionless ratio is a perfectly good rate. We grant the reply and restate the objection as a dilemma. Either the rate of passage could have been different, in which case it is measured against a second time, or it could not, in which case it is fixed by the B-facts and so cannot be what distinguishes passage from succession. Either way the rate gives no content to passage, which supports the conclusion of TSM-02 that time does not flow.
\end{abstract}

\section{The objection and the standard reply}

The objection goes back to \citet{Smart1949}. Motion is change of position with respect to time; if time itself moves, its motion is change of temporal position with respect to time, and the only rate it could have is one second per second. \citet{Markosian1993} sets out the objection and its variants and replies on behalf of passage. \citet{Maudlin2007} and \citet{Olson2009} have since disputed whether \enquote{one second per second} is a rate: Maudlin holds that it is, as a dimensionless ratio like an exchange rate of one dollar per dollar, and Olson holds that it is not, since it amounts to no more than the number one.

We do not need to decide this. Grant that one second per second is a rate. The question is then what the attribution of that rate adds to the B-facts.

\section{The dilemma}

\begin{proposition}\label{p:dilemma}
Either the rate of passage is contingent, and passage is measured against a second temporal dimension, or it is necessary, and the claim that time passes at that rate is fixed by the B-facts.
\end{proposition}

\begin{proof}
Suppose the rate is contingent: time could have passed at two seconds per second. A difference between the actual rate and the alternative would have to show up somewhere. It cannot show up in the B-series: in both cases the same events stand in the same relations of earlier and later, separated by the same durations. So it shows up only if the passage of ordinary time is timed against some further parameter, a hypertime, relative to which an ordinary second elapses in one unit or in half of one \citep{Smart1949}. Suppose instead the rate is necessary. Then \enquote{time passes at one second per second} is true in every world with a temporal metric, and in particular in every B-duplicate of ours; it is fixed by the B-facts.
\end{proof}

The first horn has a regress: if hypertime also passes, it has a rate, and so on, while if it does not, the dynamic theorist has conceded that a time can be real without passing. The second horn bears directly on the argument of TSM-02-a. By its Assumption~2.7, which the dynamic theorist accepts, passage is not fixed by the B-facts: every world has a B-duplicate without passage. A rate that every B-duplicate shares cannot be what distinguishes the world with passage from its passage-free duplicate.

\begin{corollary}\label{c:rate}
No rate of passage gives content to passage beyond succession, except at the price of hypertime.
\end{corollary}

\section{Objection and reply}

\begin{objection}
Passage is not motion through time, so it need have no rate; the objection attacks a picture no careful A-theorist holds.
\end{objection}

\begin{reply}
Then passage is a change that occurs at no rate and at no time, which is the incoherence of the A-series by another route (TSM-01). The rate objection is strongest as a challenge to say what passage adds to succession, and Corollary~\ref{c:rate} records that no rate answers the challenge. Its conclusion is the same as that of Theorem~3.2 of TSM-02-a: what is fixed by the B-facts, the impression of passage included, is no evidence that time passes.
\end{reply}

\section{Conclusion}

Whether or not one second per second is a rate, the rate of passage is either measured in hypertime or fixed by the B-facts. The first is a regress and the second is mere succession. The objection, so restated, supports the conclusion of TSM-02 that time does not flow.

\begin{thebibliography}{9}
\bibitem[Markosian(1993)]{Markosian1993} Markosian, N. (1993). How fast does time pass? \emph{Philosophy and Phenomenological Research}, 53(4), 829--844.
\bibitem[Maudlin(2007)]{Maudlin2007} Maudlin, T. (2007). \emph{The Metaphysics Within Physics}. Oxford University Press.
\bibitem[Olson(2009)]{Olson2009} Olson, E.~T. (2009). The rate of time's passage. \emph{Analysis}, 69(1), 3--9.
\bibitem[Smart(1949)]{Smart1949} Smart, J.~J.~C. (1949). The river of time. \emph{Mind}, 58(232), 483--494.
\end{thebibliography}

\end{document}
